\section{Design Goals}

The repository describes \plat{} as ``an experimental interpreter language
where code stays small, rules stay regular, and the keywords come from
Limburg.'' Its design notes emphasize simplicity, readability, composition,
and a Limburgian linguistic surface. The language is inspired by Lua, Lisp,
and Smalltalk, while aiming to remain smaller and simpler than those systems.

Several goals recur across the repository:

\begin{description}
  \item[Smallness.] The language is built from a small set of concepts:
        values, variables, expressions, functions, tables, conditionals, and
        loops. The current standard library is intentionally tiny.
  \item[Regularity.] The design avoids semicolons, complex type hierarchies,
        overloaded object models, and large feature sets. Programs execute
        top-level statements directly, and functions are global declarations.
  \item[Approachability.] The syntax is lightweight, comments use \code{\#},
        and examples are short. The implementation can print the AST with
        \code{--ast}, which supports inspection and teaching.
  \item[Cultural identity.] Core keywords use Limburgian words. The design
        notes state that the language prefers the Sittard dialect where
        possible, while recognizing that exact spelling and dialect choices
        may evolve.
  \item[Inspectability.] The implementation is organized around a parser,
        AST, interpreter, runtime values, environments, and diagnostic
        templates. This makes the artifact suitable for discussion as a small
        language implementation rather than a black-box tool.
\end{description}

The most important design distinction is between technical simplicity and
cultural intentionality. A small interpreter language with dynamic values,
lexical environments, global functions, and mutable tables is not technically
unusual. In \plat, that familiarity is intentional: it reduces the number of
semantic novelties competing for attention, allowing the Limburgian surface
syntax, diagnostics, examples, and naming choices to become the primary design
material.

% TODO: Add quotations from project maintainers or design issue discussions if
% available. Such sources would strengthen the account of project intent.
